\chapter{Measurement, calibration and the loss boundary}\label{ch:measurement}
A perfectly implemented equation can return a misleading number if its
inputs do not represent the physical event. This is why a book about code
must return to measurement: what quantity changed, where was it measured,
what uncertainty remains, and who justified the reference against which
the change is evaluated?

\section{References and scales need an interpretation}
The Gaussian reference $x_i^*$ describes a declared functional condition.
The scale $\sigma_i$ describes how deviations are weighted. Neither is
automatically an observed mean, a measurement uncertainty or a just allocation.
A physical unit conversion rescales stock, reference and scale together;
changing only the scale is a substantive change in valuation.

An empirical programme would need a parameter record explaining the
functional evidence, measurement domain, validity period and uncertainty
of each declaration. That work is not supplied by choosing convenient
numbers for a textbook example. The examples here are explicitly illustrative.

\section{Measurement and account boundaries}
Source withdrawal, useful destination delivery and loss are different
measurements. A pipeline meter at the source does not establish delivery
at the destination. A receipt should retain those distinctions even in
a first lossless model, where the declared values coincide.

Physical time matters too. An endpoint observation does not necessarily
identify the intermediate path of several simultaneous actions. If the
account uses endpoint interpolation, record that convention. If it claims
a model-realized constant-rate path, identify the physical assumption and
the evidence required to use that reading outside the model.

\section{Conservation does not choose a valuation}
An explicit loss coordinate can close total physical accounting. It does
not automatically become an argument of the homeostatic potential, and
excluding it from that potential does not automatically authorize an
additional charge. The represented consequence, process burden and physical
carrier must be reconciled without double counting.

The historical Definition 6.4 discusses a dissipated-stock charge with a
rationale that no destroyed-stock coordinate is present \cite{foundation}.
Later candidate work includes a loss ledger coordinate. The statements
``absent from state'' and ``present in a ledger but absent from $V$'' are
not literally identical. A broader loss-aware rule needs an explicit
authority/wording reconciliation; this book does not silently choose one.

The active first domain has $\eta=1$ and $C=0$. Loss-aware runtime is
not certified by that choice. An aggregate conservation residual can hide
an invalid sink update if the declared conservation profile is not enforced
in the actual transition. Carrier-specific validation remains necessary.

\section{Do not invent burdens to complete a story}
Fuel conversion, electricity delivery, pumping and water transfer can be
represented as linked atomic transformations when they are in scope. Their
existence is not a reason to add an arbitrary extra cost to every downstream
action. Each physical effect needs a declared carrier and one justified
valuation location. The chain must neither disappear inside a mega-action
nor be charged repeatedly under several names.

Monetary price, inconvenience, sanctions and physical process burden also
have different meanings. A payment may explain how an essential service
is financed without becoming a physical term in the EBU equation.

\section{Governance and privacy are not residual variables}
Selecting references, identifying owners and deciding who may access
measurements have institutional consequences. A model that calls itself
local does not automatically settle consent, privacy, auditability or access.
Those questions require an explicit governance layer and appropriate domain
review; they cannot be solved by adding a coefficient to $V$.

The practical benefit of keeping records distinct is transparency. A
necessary action can have negative represented field value without the
account declaring that its recipient should be denied service. The limited
experimental capacity policy must not be presented as a ready rule for
real-world access to essential resources.

\section*{Chapter recap}
Measurement and calibration determine what a precise calculation means.
The first lossless world allows focused mathematics while leaving, rather
than pretending to solve, the broader loss and institutional boundaries.
